\documentclass[11pt]{article}
\usepackage[a4paper, top=18mm, bottom=18mm, left=20mm, right=20mm]{geometry}
\usepackage[T2A]{fontenc}
\usepackage[utf8]{inputenc}
\usepackage[ukrainian]{babel}
\usepackage{graphicx}
\usepackage[export]{adjustbox}
\usepackage{amsmath}
\usepackage{amssymb}
\graphicspath{{/app/Quiz/Scripts/2023/ProgAll/15BinTree/}{/app/Quiz/Scripts/2021/Mod2021_3/BinTree/}{/app/Quiz/Scripts/2020/Mod2020_4/BinTree/}}
\setlength{\parindent}{0pt}
\setlength{\parskip}{6pt}
\pagestyle{empty}
\begin{document}
%begin
{\large\textbf{Контрольна робота}}

\textbf{Варіант \No\,9}\hfill \textit{ПІБ:}\,\underline{\hspace{60mm}}
\vspace{4mm}

%beginitem
1. Що буде виведено за виконання фрагменту коду?

%code
print('R', end=' ')
n=3
while n<7:
    n+=1
print(n)
%code

\vspace{5mm}
%enditem

%beginitem
2. Що буде виведено за виконання фрагменту коду?

%code
n=3
while n<7:
    n+=1
print(n)
%code

\vspace{5mm}
%enditem

%beginitem
3. %goodpath:Scripts/2018/Mod2018_1/03-good.txt
Відмітити все, що є коректним оператором С++:
( 1 ) cout<<(3=!5);

( 2 ) bool f(int a);

( 3 ) cin<<x;

( 4 ) j=i;
\vspace{5mm}
%enditem

%beginitem
4. Що буде виведено за виконання фрагменту коду?
code
_d = 0

class D:
    _d = 1
    
    def __init__(self):
        self._d = 2
        _d = 3
        D._d = 4

obj = D()
try:
    print(_d, end=' ')
    print(obj._d, end=' ')
    print(D._d, end=' ')
    print(obj._d, end=' ')
    print(obj._D__d, end=' ')
except:
    print('ERR')
code
\vspace{5mm}
%enditem

%beginitem
5. %goodpath:Scripts/2019/Mod2019_1/Lex/03-good.txt
Відмітити все, що є коректним оператором С++:
( 1 ) if a<b: a=0 else: b=0

( 2 ) x=3**1

( 3 ) if ('x'><'y'): a=6

( 4 ) --n
\vspace{5mm}
%enditem

%beginitem
6. Що буде виведено за виконання фрагменту коду?

%code
print('R', end=' ')
d=[10,11,12,13,14]
del d[-6:-6]
print(d)
%code

\vspace{5mm}
%enditem

%beginitem
7. %goodpath:Scripts/2019/Mod2019_1/Lex/01-good.txt
Відмітити все, що є однією правильною лексемою:
( 1 ) (1)

( 2 ) 0.2E3

( 3 ) "abc'

( 4 ) x23
\vspace{5mm}
%enditem

%beginitem
8. Що буде виведено за виконання фрагменту коду?

%code
print('R', end=' ')
n=7
while n>3:
    n+=-3
print(n)
%code

\vspace{5mm}
%enditem

%beginitem
9. Що буде виведено за виконання фрагменту коду?
%code
if (3 >= 3)
    cout << "x";
else
    cout << "b";
%code

\vspace{5mm}
%enditem

%beginitem
10. Що буде виведено за виконання фрагменту коду?

%code
def h(a,b):
    c=69
    if b>=-3:
        return 3
    if a<=-4:
         c=1
    else: 
        return 5
    return c

print(h(9,9))
%code
\vspace{5mm}
%enditem

%end
\newpage
\end{document}

